\documentclass{article}
\usepackage[utf8]{inputenc}
\usepackage[T1]{fontenc}
\usepackage{geometry}
\geometry{
    a4paper,
    left=2.5cm,
    right=2.5cm,
    top=2.5cm,
    bottom=2.5cm
}
\usepackage{graphicx}
\usepackage{float}
\usepackage{hyperref}
\usepackage[english]{babel}
\usepackage{amsmath}
\usepackage{amssymb}
\usepackage{array}
\usepackage{booktabs}
\usepackage{caption}
\usepackage{subcaption}
\usepackage{enumitem}
\usepackage{fancyhdr}
\usepackage{titlesec}
\usepackage{url}
\usepackage{textcomp}
\usepackage{xcolor}
\usepackage{listings}
\usepackage{multirow}
\usepackage{longtable}
\usepackage{tabularx}
\usepackage{times}

% Set up page style
\pagestyle{fancy}
\fancyhf{}
\rfoot{\thepage}
\renewcommand{\headrulewidth}{0pt}

% Title formatting
\titleformat{\chapter}[display]{\normalfont\bfseries}{\chaptertitle\thechapter}{0pt}{\Huge}
\titlespacing*{\chapter}{0pt}{50pt}{40pt}

% Document info
\title{Template Document}
\author{Academic Support Agent}
\date{\today}

\begin{document}
\maketitle

\begin{center}

\includegraphics[width=2.17cm,height=2.14cm]{ufcg.png}


{\textbf {UNIVERSIDADE FEDERAL DE CAMPINA GRANDE\\
PR -REITORIA DE PESQUISA E EXTENS O\\
COORDENA O DE PESQUISA}}
\end{center}

 
\vspace{24pt}
\begin{mybox}
\begin{center}
{\Large \textbf{RELAT RIO DE ATIVIDADES DO ALUNO}}
\end{center}
\end{mybox}
\vspace{12pt}
%%%%%%%%%%%%%%%%%%%%%%%%%%%%%%%%%%%%%%%%%%
%Informa es do Projeto
%%%%%%%%%%%%%%%%%%%%%%%%%%%%%%%%%%%%%%%%%%


\begin{flushleft}
{\large \textbf{Programa:} PIBIC}\\
\vspace{14pt}
{\large \textbf{T tulo do Projeto:} Meu Projeto}\\
\vspace{14pt}
{\large \textbf{Aluno(a):} Nome do(a) Aluno(a)}\\
\vspace{14pt}
{\large \textbf{Orientador(a):} Nome do(a) orientador(a) }
\end{flushleft}
\vspace{200pt}

%%%%%%%%%%%%%%%%%%%%%%%%%%%%%%%%%%%%%%
%Informa es Assinaturas
%%%%%%%%%%%%%%%%%%%%%%%%%%%%%%%%%%%%%%
\signature{Sr(a). Bolsista}{Bolsista}
\signature{Dr(a). Orientador(a)}{Orientador(a)}

\insertdate[Sum ]


\newpage

\pagenumbering{gobble}

\section*{Introdu o}

 Em~\citeonline{RCT2017}
existe uma cole o de manuais completos sobre quest es relacionadas a instala o e utiliza o do \textbf{R}, ver em \url{https://cran.r-project.org/manuals.html}.   poss vel encontrar v rios livros com aplica es da linguagem \textbf{R} em diferente  rea, por exemplo, ver~\citeonline{Paradis2004,Baddeley2004, Morandat2012}.
Pacotes \textbf{R} s o compostos de bibliotecas com fun es e dados.


\section*{Objetivos}

\section*{Material e M todos/Metodologia}

\section*{Resultados Parciais}

\newpage
\end{document}